%%
\section{Thematic Literature Analysis}
\label{sec:thematic-analysis}
%%

\autoref{tab:tracking-topics} shows the distribution of papers across 40 topics generated during the thematic analysis of the web tracking literature from 2005 to 2024, published at one of the seven top security and privacy venues--IEEE S\&P, USENIX Security, ACM CCS, NDSS, ACM IMC, PETS, and WWW. This thematic organization is available at \href{https://osf.io/zsy7e/overview?view_only=a346b98bf89244b5b17d842f6fb3fb13}{this URL}.

\begin{table}
\centering
\small
\renewcommand{\arraystretch}{1.15}

\begin{tabular}{l r}
\hline
\noalign{\vskip 1.5pt}
\textbf{Topic} & \textbf{\# Papers} \\
\noalign{\vskip 1.5pt}
\hline
\noalign{\vskip 1.5pt}
Tracking Measurements & 50 \\
Tracking in Browsers & 23 \\
Browser Fingerprinting Measurement/Detection & 20 \\
Regulation Compliance & 20 \\
Online Profiling & 19 \\
Tracking in Mobile & 19 \\
Third-party Tracking & 18 \\
Cookie Consent/Banners & 17 \\
Cookie Measurements & 17 \\
User Studies on Tracking & 15 \\
Ad Blocking Measurements & 13 \\
Javascript-based Tracking & 10 \\
Side Channel Tracking & 10 \\
ATS Defense Solutions & 9 \\
Browser Extension Fingerprinting & 8 \\
Tracker-Advertiser Flows & 7 \\
URL-based Tracking & 7 \\
Browser Fingerprinting Defense & 6 \\
Cross-site Tracking & 6 \\
Anti Ad Blocking & 5 \\
Same-site Tracking & 5 \\
Tracking Analysis Tools & 5 \\
Cross-device Tracking & 4 \\
Device Fingerprinting & 4 \\
PII-based Tracking & 4 \\
Privacy Sandbox & 4 \\
Ad Blocking Solution & 3 \\
Bounce Tracking & 3 \\
Breakage Studies & 3 \\
First-party Tracking & 3 \\
Privacy Policy Analysis & 3 \\
Privacy Studies of Browser Extensions & 3 \\
Search Engine Privacy & 3 \\
Behavioral Tracking & 2 \\
Cookie Blocking & 2 \\
Cookie Syncing & 2 \\
DNS Tracking & 2 \\
Form-based Tracking & 2 \\
Location-based Tracking & 2 \\
Pixel Tracking & 2 \\
\noalign{\vskip 1.5pt}
\hline
\end{tabular}

\caption{Distribution of papers across different topics.}
\label{tab:tracking-topics}
\end{table}


%%
\section{Measurement Methodologies}
\label{sec:measurement-methodologies}

%%
\subsection{Crawling Measurements}
\label{sec:crawling}

Web crawling with instrumented browsers is the most common approach to measure online tracking.
Browser instrumentation can take two forms: \textit{out-of-band} or \textit{in-band}. Out-of-band, or deep instrumentation, modifies directly the browser or JavaScript engine. In contrast, in-band leverages instrumentation hooks, like prototype patching, at the JavaScript level to overwrite functionality of interest.

\para{User Agent.} Most measurements require a browser supporting modern web features.
Simplified user agents which do not execute JavaScript or have incomplete support for web APIs can be appropriate for targeted measurements (\eg{}, Apache Nutch\footnote{\url{https://nutch.apache.org/}}). 

\para{Automation Frameworks with Instrumentation Hooks.} To drive full consumer browsers, researchers rely on automation tooling built for website and browser testing: \eg{}, Chrome DevTools Protocol (CDP) for Blink-based browsers and Marionette for Gecko-based browsers. These internal interfaces are used by cross-browser automation libraries like Selenium or Puppeteer~\cite{
cross-browser-testing-1-2020,puppeteer-support-firefox-2024}.
Many researchers make direct use of these libraries, while several projects which bundle full browser automation with additional instrumentation and measurement tooling also exist~\cite{englehardtOnlineTracking1millionsite2016,openwpm4StudiesUsing,mayerFourthpartyFourthpartyFourthParty2011}.


\para{Deep Instrumentation.} Many attempts have been made in leveraging deep instrumentation for security-related web measurements~\cite{neasbittWebCapsuleLightweightForensic2015,liJSgraphEnablingReconstruction2018, acarFPDetectiveDustingWeb2013, IncontextStoreIncontext1998, chenMystiqueUncoveringInformation2018}. The fundamental problem is that the browser evolves rapidly, rendering research prototypes obsolete quickly, as maintaining the patches is difficult or impossible~\cite{chenMystiqueUncoveringInformation2018}.  
Two major efforts try to overcome this limitation: VisibleV8~\cite{jueckstockVisibleV8InbrowserMonitoring2019} and PageGraph\footnote{\url{https://github.com/brave/brave-browser/wiki/PageGraph}}.
%
PageGraph is maintained directly by the Brave Browser team, making it the only deep instrumentation framework that has browser support. 
%
VisibleV8\footnote{\url{https://github.com/wspr-ncsu/visiblev8}} is designed so that its patches are minimal (67 lines of code for the actual JavaScript monitoring) and has been successful in providing builds from Chromium 63 to 137 (version at submission time) with minimal effort~\cite{ jueckstockVisibleV8InbrowserMonitoring2019}. A major benefit is that deep instrumentation is agnostic to what needs to be monitored: \textit{all} web APIs can be hooked, even when not knowing the responsible APIs beforehand~\cite{suAutomaticDiscoveryEmerging2023}.

\para{Stealthiness.}
A significant threat to the validity of active web measurements is the ability of websites to detect crawlers and instrumented browsers. Upon detection, websites may block crawlers or alter their behavior (a practice known as \textit{cloaking})~\cite{invernizzi2016cloak}. Automation frameworks often inject detectable artifacts in the JavaScript context or alter the user-agent string. Researchers may need to deploy further evasion techniques\footnote{\url{https://www.npmjs.com/package/puppeteer-extra-plugin-stealth}} to avoid differential treatment.

\para{Site lists.} Top lists of popular websites are published by several sources based on different methodologies: Alexa Top Million (now deprecated), Cisco Umbrella Popularity List, Majestic Million, Tranco~\cite{lepochatTrancoResearchOrientedTop2019}, Google CrUX, or Cloudflare Radar. 
%
Their use as a proxy to study websites and real users' behaviors has raised some skepticism in the past as these lists can be unstable, inconsistent, and prone to manipulation. Moreover, the choice of top list can sometimes impact research findings~\cite{lepochatTrancoResearchOrientedTop2019,ruthTopplingTopLists2022}.

\para{Existing Crawl Datasets.}
Another strategy is to leverage existing web crawl datasets. Nonprofit organizations and community-driven projects such as the Internet Archive\footnote{\url{https://archive.org/}}, Common Crawl\footnote{\url{https://commoncrawl.org/}}, and the HTTP Archive\footnote{\url{https://httparchive.org/}} routinely crawl websites and publish their data openly. Other initiatives such as \textit{privacytests.org}\footnote{\url{https://privacytests.org}} perform continuous open testing of web browser privacy across different vendors and tracking vectors, but do not include more sophisticated tracking approaches. 

\para{Limitations.} Representativeness and generalizability issues arise due to 
bot detection measures~\cite{krumnowHowGullibleAre2022},
measurement vantage points~\cite{samarasingheGlobalPerspectiveWeb2019},
device form factors~\cite{yangComparativeMeasurementStudy2020,casselOmniCrawlComprehensiveMeasurement2022}
and potential differences in results obtained from crawls versus real browsing by humans~\cite{zeberRepresentativenessAutomatedWeb2020}. Directly related, studies are often very difficult to reproduce and replicate as differences in methodologies and experimental setups are not always fully documented by researchers~\cite{demirReproducibilityReplicabilityWeb2022,hantkeWebExecutionBundles2025}.


%%
\subsection{User Studies}
\label{sec:user-studies}
%%

In practice, user studies can take multiple forms; they can be conducted through \textit{usability surveys or interviews}, or be based on data collected from real users through \textit{field measurements, crowdsourcing, or direct collection} through a browser extension or application. As an example, the National Internet Observatory\footnote{\url{https://nationalinternetobservatory.org/index.html}}, a nascent effort, invites US residents to volunteer data about their online behaviors and allows privacy-preserving access to researchers for scientific studies~\cite{steningUnprecedentedDataCollection2022,callahanCanWeBetter2021,fealIntroductionNationalInternet2024}.
%
With these techniques, researchers have mostly investigated participants’ comprehension, perception, and interaction with respect to cookie  dialogs~\cite{birrellSoKTechnicalImplementation2024,machuletzMultiplePurposesMultiple2020,bermejofernandezThisWebsiteUses2021,habibOkayWhateverEvaluation2022,singhWhatCookieConsent2022,bielovaSurveyAcademicStudies2022,Bielova2024-zr}. They typically study and compare different consent dialog designs, finding that many current designs effectively nudge participants towards more privacy-preserving options~\cite{machuletzMultiplePurposesMultiple2020,bermejofernandezThisWebsiteUses2021}. These studies also recommend that consent choices be \textit{reject by default} and that users should be able to easily revisit choices they have made~\cite{habibOkayWhateverEvaluation2022,Kanc-etal-25-PETs}.


\section{Need for Fingerprinting}
\label{sec:need-for-fingerprinting}

%%
Constructively, fingerprinting can be thought of as a form of intrusion detection. 
%
Web applications can learn about the browsing environment or behavior of their first-party users and associate it with specific user identities~\cite{LinPhishSheepsClothing2022}. 
%
For example, a website can learn that the user is using a smartphone with specific dimensions or a desktop browser with a specific kind of GPU. 
%
If the user's credentials are ever stolen and an attacker attempts to login to that service, the service can extract the attacker's fingerprint, observe a major difference against fingerprints from prior sessions, and request additional authentication data from the attacker (such as a one-time password).
%
The same techniques can be constructively used to differentiate real users from malicious bots or attackers engaging in ad fraud.
%
In practice, large online services commonly incorporate fingerprinting-derived signals into broader abuse-mitigation pipelines to detect automated account creation, credential stuffing, and spam. 
%
These signals are typically combined with behavioral features and rate-limiting mechanisms, enabling services to adapt authentication requirements based on assessed risk and maintaining service integrity in adversarial environments where traditional identifiers are unreliable.

%%
Destructively, the same techniques that can identify user-impersonating attackers and bots, turning against users who wish to keep their identity anonymous. 
%
With fingerprinting, a web app may be able to determine that a certain anonymous user is in fact eponymous, since their browser fingerprint matches that of a known user on the same platform. 
%
This undesired re-identification occurs despite user's attempt to hide by deleting cookies or using the browser's private mode. 
%
In a cross-site context, fingerprinting can be abused to link unrelated website visits, even when 3P cookies are disabled.


% ---------- Preamble helpers (once) ----------
% \usepackage{booktabs}
% \usepackage{tabularx}
% \usepackage{array}
\newcolumntype{Y}{>{\raggedright\arraybackslash}X}

% =========================================================
% Appendix Table A1: Stateful / identifier-based tracking
% =========================================================
\begin{table*}[t]
\centering
\scriptsize
\setlength{\tabcolsep}{3pt}
\renewcommand{\arraystretch}{1.15}
\begin{tabularx}{\textwidth}{p{2.8cm} Y Y}
\toprule
\textbf{Tracking Mechanism} & \textbf{How does it work? / What does it enable?} & \textbf{Mitigations / Defenses} \\
\midrule

\multicolumn{3}{l}{\textbf{Third-party stateful tracking}}\\
\addlinespace[2pt]

Third-party cookies &
Cross-site identifiers set/read by embedded third parties enable recognition across unrelated sites and support profiling/retargeting at scale \cite{roesnerDetectingDefendingThirdparty2012,englehardtOnlineTracking1millionsite2016,gomerNetworkAnalysisThird2013,frankenWhoLeftOpen2018,beugin2024WebAlmanac2024}. &
Block or partition third-party cookies and other state; use interaction-gated access (e.g., Storage Access API), enforce cookie attributes (e.g., SameSite), and deploy tracking protections/storage partitioning \cite{TrackingPreventionWebKit2020,privacycommunitygroupClientSideStoragePartitioning2022,mdnStorageAccessAPI2024,khodayariStateSameSiteStudying2022}. \\

Cookie syncing (ID syncing) &
Trackers/ad exchanges reconcile distinct user IDs by sharing identifiers via redirects, URL parameters, or back-channel coordination, expanding the effective tracking graph \cite{papadopoulosCookieSynchronizationEverything2019,bashirTracingInformationFlows2016}. &
Partition/block third-party state; detect and block syncing endpoints and redirect patterns; use flow/graph-based approaches to identify identifier-sharing structures \cite{sibyWebGraphCapturingAdvertising2022,iqbalAdGraphGraphBasedApproach2020}. \\

Tracking pixels &
Tiny/invisible resources trigger requests that transmit identifiers and page context; widely used for measurement and ad attribution \cite{yuTrackingTrackers2016,fouadMissedFilterLists2020}. &
Endpoint and content blocking (plus better filter maintenance); URL/parameter sanitization to reduce identifier leakage in requests and navigations \cite{fouadMissedFilterLists2020,munirPURLSafeEffective2024,bekosHitchhikersGuideFacebook2023}. \\

Session replay &
Third-party JS records DOM events and user interactions (including form inputs) for analytics/support; misconfiguration can capture sensitive data \cite{englehardtNoBoundariesCredentials2018}. &
Block or constrain session-replay vendors and scripts; limit access to sensitive fields; auditing/runtime controls and policy enforcement to prevent over-collection \cite{englehardtNoBoundariesCredentials2018}. \\

\addlinespace[3pt]
\multicolumn{3}{l}{\textbf{First-party stateful tracking}}\\
\addlinespace[2pt]

% (respawning)
Supercookies / evercookies &
Identifiers are redundantly stored across multiple APIs (e.g., Flash/legacy storage, localStorage, caches/ETags, favicon/cache tricks, HSTS-like bits) to recreate IDs after deletion \cite{kamkarSamyKamkarEvercookie2010,ayensonFlashCookiesPrivacy2011,fouadMyCookiePhoenix2022,solomosTalesFaviconsCaches2021,mishraDejaVuAbusing2021,syversonHSTSSupportsTargeted2018}. &
Comprehensive state clearing beyond cookies; deprecate legacy plugins; partition client-side storage/network state; use standards like storage partitioning and “clear site data” style policies \cite{privacycommunitygroupClientSideStoragePartitioning2022,TrackingPreventionWebKit2020}. \\

First-party cookies &
Embedded third-party scripts write IDs into the \emph{first-party} cookie jar (or use first-party endpoints), bypassing third-party cookie blocking and enabling site-scoped identifiers \cite{chenCookieSwapParty2021,demirUnderstandingFirstPartyCookie2022}. &
Attribute cookies to responsible scripts and detect tracking-purpose cookies; isolate/partition or restrict first-party storage writes by third-party code \cite{munirCookieGraphUnderstandingDetecting2023,bahramiCOOKIEGUARDCharacterizingIsolating2024}. \\

CNAME cloaking &
A first-party subdomain is CNAMEd to a third-party tracker so requests/storage appear first-party, evading naive third-party heuristics \cite{daoCNAMECloakingBasedTracking2021,dimovaCNAMEGameLargescale2021}. &
Resolve CNAME chains and treat cloaked services as third-party; apply list + behavior-based detection and partition/block accordingly \cite{daoCNAMECloakingBasedTracking2021,dimovaCNAMEGameLargescale2021}. \\

Link decoration &
Identifiers or join keys are appended to URLs and propagated across navigations/redirects, enabling linkage across sites and contexts \cite{munirPURLSafeEffective2024,bekosHitchhikersGuideFacebook2023}. &
Sanitize tracking parameters; constrain identifier propagation through redirects and navigations; use allow/deny lists + heuristics for tracking params \cite{munirPURLSafeEffective2024}. \\

Bounce tracking &
Intermediate “bounce” domains perform redirects while reading/writing state, enabling cross-site linkage despite stricter cookie controls \cite{braveprivacyteamUnlinkableBouncingMore2022,kellyBounceTrackingMitigations2022}. &
Navigation tracking mitigations that detect bounces and purge/partition storage for bounce domains; limit cross-site state access around bounce chains \cite{braveprivacyteamUnlinkableBouncingMore2022,kellyBounceTrackingMitigations2022}. \\

Server-side tracking &
Event reporting shifts from client-side pixels to server-to-server pipelines, reducing the visibility/effectiveness of browser-side blockers \cite{fouadDevilDetailsDetection2024,fraihiClientsideServersideTracking2024}. &
Detect server-side endpoints and event pipelines via measurement; combine technical auditing with stronger governance/compliance controls \cite{fouadDevilDetailsDetection2024,fraihiClientsideServersideTracking2024}. \\

\bottomrule
\end{tabularx}
\caption{\textbf{Stateful / identifier-based} tracking mechanisms and common mitigations.}
\label{tab:appendix-stateful-tracking}
\end{table*}

% =========================================================
% Appendix Table A2: Stateless tracking (fingerprinting, sensors, side-channels)
% =========================================================
\begin{table*}[t]
\centering
\scriptsize
\setlength{\tabcolsep}{3pt}
\renewcommand{\arraystretch}{1.15}
\begin{tabularx}{\textwidth}{p{2.8cm} Y Y}
\toprule
\textbf{Tracking Mechanism} & \textbf{How does it work? / What does it enable?} & \textbf{Mitigations / Defenses} \\
\midrule

Browser fingerprinting &
Derives stable identifiers from exposed attributes/APIs (UA, fonts, canvas, WebGL, audio, timing, etc.), often without any stored client state \cite{effpanopticlick,nikiforakisCookielessMonsterExploring2013,acarFPDetectiveDustingWeb2013,laperdrixBeautyBeastDiverting2016,vastelFPSTALKERTrackingBrowser2018}. &
Reduce entropy (standardize outputs, clamp timers, permission-gate APIs); anti-fingerprinting (“lie/perturb”) and privacy-budget/noise approaches \cite{gomez-boixHidingCrowdAnalysis2018,dattaEvaluatingAntiFingerprintingPrivacy2019,dotyMitigatingBrowserFingerprinting2019,nikiforakisPriVaricatorDeceivingFingerprinters2015,snyderBraveFingerprintingPrivacy2019,torANTIFINGERPRINTINGTorProject,firefoxResistFingerprinting,rescorlaTechnicalCommentsPrivacy2021}. \\

Canvas fingerprinting &
Uses pixel readback from rendered canvases to infer device/software differences, frequently combined with other signals \cite{moweryPixelPerfectFingerprinting2012,acarFPDetectiveDustingWeb2013}. &
Restrict readback, standardize rendering, or add noise; deploy anti-fingerprinting modes \cite{dattaEvaluatingAntiFingerprintingPrivacy2019,torANTIFINGERPRINTINGTorProject}. \\

WebGL / GPU fingerprinting &
Extracts high-entropy signals from GPU/driver rendering behavior (e.g., shader outputs), enabling robust identifiers \cite{wuRenderedPrivateMaking2019,laorDRAWNAPARTDevice2022}. &
Uniform rendering pipelines, reduce surface area, clamp high-resolution timing, and apply privacy budgets/noise \cite{rescorlaTechnicalCommentsPrivacy2021,dotyMitigatingBrowserFingerprinting2019}. \\

Audio fingerprinting &
Captures distinctive characteristics from the audio stack/hardware pipeline to identify devices \cite{chaliseYourSpeakerMy2022}. &
Permission-gate and standardize/noise audio outputs; limit high-entropy API exposure \cite{rescorlaTechnicalCommentsPrivacy2021}. \\

% (motion/IMU)
Sensor-based fingerprinting &
Uses accelerometer/gyroscope access to infer device properties or behavioral signals that aid tracking \cite{dasTrackingMobileWeb2016,dasWebsSixthSense2018}. &
Reduce precision and sampling rates; require permissions and add noise \cite{dasWebsSixthSense2018}. \\

Battery status leakage &
Battery APIs expose fine-grained state that can form short-term identifiers and facilitate linkage \cite{olejnikLeakingBatteryPrivacy2016,olejnikBatteryStatusNot2017}. &
Coarsen/remove battery APIs; permission-gate and rate-limit access \cite{olejnikBatteryStatusNot2017}. \\

Timing / cache side channels &
Timing/cache effects can reveal cross-origin state or contribute to fingerprinting even under some restrictions \cite{sanchez-rolaClockClockTimeBased2018,weinbergStillKnowWhat2011,shusterman2021prime+,jinTimingBasedBrowsingPrivacy2022}. &
Timer clamping/jitter, stronger isolation, and reducing shared resources to limit leakage \cite{sanchez-rolaClockClockTimeBased2018,rescorlaTechnicalCommentsPrivacy2021}. \\

Extension fingerprinting &
Detects installed extensions via web-accessible resources and behavioral probes, yielding high-entropy identifiers \cite{starovXHOUNDQuantifyingFingerprintability2017,trickelEveryoneDifferentClientside2019,karamiCarnusExploringPrivacy2020,karamiUnleashSimulacrumShifting2022,agarwalPeekingWindowFingerprinting2024}. &
Harden extension resource policies; reduce observable surfaces and isolate extension contexts \cite{laperdrixFingerprintingStyleDetecting2021}. \\

Behavioral fingerprinting &
Builds identifiers from user interaction dynamics (mouse movements, typing patterns), often complementing other signals \cite{leivaMyMouseMy2021,yuanCrossdeviceTrackingIdentification2018}. &
Limit granularity and frequency of event streams; detect/block behavioral scripts; minimize exposure of high-frequency events \cite{leivaMyMouseMy2021}. \\

Cross-device tracking &
Links identities across devices using deterministic identifiers (logins/emails) and/or probabilistic matching of signals and behaviors \cite{diaz-moralesCrossDeviceTrackingMatching2015,brookmanCrossDeviceTrackingMeasurement2017,zimmeckPrivacyAnalysisCrossdevice2017,wangGraphTrackGraphbasedCrossDevice2022}. &
Reduce reuse of stable join keys; constrain identifier sharing; strengthen transparency/controls and governance \cite{brookmanCrossDeviceTrackingMeasurement2017,zimmeckPrivacyAnalysisCrossdevice2017}. \\

\bottomrule
\end{tabularx}
\caption{\textbf{Stateless} tracking mechanisms (fingerprinting, sensors, side-channels) and common mitigations.}
\label{tab:appendix-stateless-tracking}
\end{table*}

% % =========================================================
% % Appendix Table A3: Detection, blocking, and compliance mechanisms
% % =========================================================
% \begin{table*}[t]
% \centering
% \scriptsize
% \setlength{\tabcolsep}{3pt}
% \renewcommand{\arraystretch}{1.15}
% \begin{tabularx}{\textwidth}{p{2.8cm} Y Y}
% \toprule
% \textbf{Theme} & \textbf{Tracking-related challenge / risk (with citations)} & \textbf{Mitigations / approaches (with citations)} \\
% \midrule

% Ecosystem measurement &
% Tracking is multi-party and dynamic (redirect chains, inclusion graphs, shifting endpoints), complicating attribution and policy enforcement \cite{englehardtOnlineTracking1millionsite2016,lernerInternetJonesRaiders2016,yuTrackingTrackers2016}. &
% Large-scale measurement and improved instrumentation to characterize behaviors and flows in situ \cite{englehardtOnlineTracking1millionsite2016,yuTrackingTrackers2016}. \\

% Blocklists and evasion &
% Filter lists can under-block and over-block; trackers adapt via domain churn and “first-party cover” techniques \cite{iqbalAdWarsRetrospective2017,fouadMissedFilterLists2020}. &
% Maintain lists and detect gaps \cite{fouadMissedFilterLists2020}; complement static rules with automated generation and behavioral methods \cite{leAutoFRAutomatedFilter2023}. \\

% Graph/flow-based blocking &
% URL heuristics miss identifier sharing and indirect data flows across the ad/analytics graph \cite{bashirTracingInformationFlows2016,papadopoulosCookieSynchronizationEverything2019}. &
% Graph-/flow-based attribution and blocking to capture structure beyond URLs \cite{iqbalAdGraphGraphBasedApproach2020,sibyWebGraphCapturingAdvertising2022,yangWTAGRAPHWebTracking2022}. \\

% JS blocking vs. breakage &
% Aggressive JS blocking may reduce tracking but can cause significant site breakage and UX regressions \cite{amjadBlockingJavaScriptBreaking2023}. &
% Selective blocking and replacement strategies; evaluate breakage systematically to improve deployability \cite{amjadBlockingJavaScriptBreaking2023}. \\

% Cookie consent UX and dark patterns &
% Consent dialogs can be manipulative and non-compliant; users misunderstand controls; revocation and “do not sell” flows are inconsistent \cite{dabrowskiMeasuringCookiesWeb2019,sanchezrolaCanIOptOutYet2019,utzUninformedConsentStudying2019,habibOkayWhateverEvaluation2022,rasaiiThouShaltNot2023,khandelwalAutomatedCookieNotice2023,sorensenGDPRChangesThird2019,vannortwickSettingBarLow2022,carpinetoAutomaticAssessmentWebsite2016,ouViopolicyDetectorAutomatedApproach2022}. &
% Measure compliance and detect dark patterns; improve UX and enforcement via automated auditing and clearer requirements \cite{habibOkayWhateverEvaluation2022,rasaiiThouShaltNot2023,khandelwalAutomatedCookieNotice2023}. \\

% Global Privacy Control (GPC) &
% A standardized opt-out signal must be usable and enforceable; deployments and compliance vary \cite{berjonGPCGDPR2021,zimmeckUsabilityEnforceabilityGlobal2023,zimmeckGlobalPrivacyControl2024}. &
% Design and evaluate opt-out signals and enforcement mechanisms; build generalizable active privacy control frameworks \cite{zimmeckGeneralizableActivePrivacy2024,zimmeckUsabilityEnforceabilityGlobal2023,zimmeckGlobalPrivacyControl2024}. \\

% User-facing cookie controls &
% Users want to block tracking cookies without breaking websites, but mental models and UI design matter \cite{schoniBlockCookiesNot2024}. &
% Design/evaluate cookie-control tools and improve usability of privacy controls \cite{schoniBlockCookiesNot2024}. \\

% Arms-race dynamics &
% As browsers restrict third-party state, trackers shift to first-party techniques, link decoration, server-side tracking, and fingerprinting \cite{narayanan2018web,callahanCanWeBetter2021}. &
% Holistic mitigations combining technical controls, measurement, and governance to anticipate adaptation pathways \cite{callahanCanWeBetter2021,narayanan2018web}. \\

% \bottomrule
% \end{tabularx}
% \caption{Appendix A3: Literature on \textbf{measurement, blocking, and compliance} intersecting with web tracking and mitigations (citations embedded inline).}
% \label{tab:appendix-mitigation-literature}
% \end{table*}
